\par\Needspace{6\baselineskip}\medskip
\begin{center}\hypertarget{native-230-article167}{}\pdfbookmark[2]{167.}{article-167}167.\end{center}
Der Darstellung der Lösung stellen wir eine Liste der wichtigsten Zeichen voran. Wir bezeichnen mit
\par\noindent
\begin{tabular}{@{}ll@{}}
$t,\ t',\ t'',\ t'''$ & die Zeitpunkte der vier Beobachtungen\\
$a,\ a',\ a'',\ a'''$ & die geozentrischen Längen des Himmelskörpers\\
$\beta,\ \beta',\ \beta'',\ \beta'''$ & seine Breiten\\
$r,\ r',\ r'',\ r'''$ & die Abstände von der Sonne\\
$\rho,\ \rho',\ \rho'',\ \rho'''$ & die Abstände von der Erde\\
$l,\ l',\ l'',\ l'''$ & die heliozentrischen Längen der Erde\\
$B,\ B',\ B'',\ B'''$ & die heliozentrischen Breiten der Erde\\
$R,\ R',\ R'',\ R'''$ & die Abstände der Erde von der Sonne
\end{tabular}
\par\noindent\SourceLeaf{231}
$(n01),\ (n12),\ (n23),\ (n02),\ (n13)$ die doppelten Flächeninhalte der Dreiecke, die jeweils
die Sonne und den ersten und zweiten,
zweiten und dritten, dritten und vierten, ersten und dritten, zweiten und vierten Ort des Himmelskörpers als Eckpunkte haben
\par\noindent
$\dfrac{1}{(\eta01)}\mathpunct{\raisebox{.45ex}{,}}\quad
\dfrac{1}{(\eta12)}\mathpunct{\raisebox{.45ex}{,}}\quad
\dfrac{1}{(\eta23)}$ die Quotienten, die sich aus der Division der Flächeninhalte
\newline
$\frac12(n01),\ \frac12(n12),\ \frac12(n23)$
durch die Flächeninhalte der entsprechenden Sektoren ergeben
\[
P'=\frac{(n12)}{(n01)}\mathpunct{\raisebox{.45ex}{,}}\qquad
P''=\frac{(n12)}{(n23)}
\]
\[
\begin{aligned}
Q'&=\left(\frac{(n01)+(n12)}{(n02)}-1\right)r'^3,\\[.5em]
Q''&=\left(\frac{(n12)+(n23)}{(n13)}-1\right)r''^3
\end{aligned}
\]
\noindent $v,\ v',\ v'',\ v'''$ die Längen des Himmelskörpers auf seiner Bahn, von einem beliebigen Punkt aus gezählt.
\par
Schließlich bezeichnen wir für die zweite und dritte Beobachtung die heliozentrischen Orte der Erde auf
der Himmelskugel mit $A'$, $A''$, die geozentrischen Orte des Himmelskörpers
mit $B'$, $B''$ und seine heliozentrischen Orte mit $C'$, $C''$.
\par
Mit diesen Bezeichnungen geht es wie bei der Aufgabe des vorigen Abschnitts
(Artikel 136) zunächst darum, die Lage der Großkreise $A'C'B'$,
$A''C''B''$ zu bestimmen. Ihre Neigungen gegen die Ekliptik bezeichnen wir mit $\gamma'$, $\gamma''$. In derselben
Rechnung bestimmen wir auch die Bögen $A'B'=\delta'$, $A''B''=\delta''$. Daraus
folgt unmittelbar
\[
\begin{aligned}
r'&=\surd(\rho'\rho'+2\rho'R'\cos\delta'+R'R')\\
r''&=\surd(\rho''\rho''+2\rho''R''\cos\delta''+R''R''),
\end{aligned}
\]
\noindent Oder wir setzen $\rho'+R'\cos\delta'=x'$, $\rho''+R''\cos\delta''=x''$,
$R'\sin\delta'=a'$, $R''\sin\delta''=a''$ und erhalten
\[
\begin{aligned}
r'&=\surd(x'x'+a'a')\\
r''&=\surd(x''x''+a''a'').
\end{aligned}
\]
